\section{The two infinite-family applications}
\subsection{Chiral square torus maps}
For the chiral maps $\{4,4\}_{(b,c)}$ in the usual square-torus parameter convention, the source assumes $bc(b-c)\ne0$. We retain the standard identification with a quotient of the square tessellation. In that identification the oriented marked group is
\[
     \bigl(\mathbb Z[i]/(b+ic)\bigr)\rtimes\langle i\rangle,
 \qquad u:z\mapsto iz,\qquad v:z\mapsto iz+1.
\]
The affine notation describes elements of the abstract semidirect product; even a degenerate small quotient is not being declared an abstract polytope merely from its permutation action. The product $vu^{-1}$ is translation by $1$, and conjugation by $u$ gives translation by $i$, so the indicated elements generate all translations and the quarter-turn. Their products satisfy $(uv)^2=1$. Also $u^4=v^4=1$, since the successive translation terms for $v^4$ sum to $1+i-1-i=0$. The linear-part homomorphism $\ell$ onto $C_4$ sends both $u$ and $v$ to $i$, so both generators have exact order four.

On the universal marked group the mirror replaces $(u,v)$ by $(u^{-1},u^2v)$. Both new linear parts are $i^{-1}$. Thus $\ell\mu$ is inversion composed with $\ell$. The relation kernel and its mirror both lie in the kernel of this linear-part map. Their product does too, and consequently
\[
                         X(K)\leq\ker\ell.
\]
But $\ell$ is injective on $\langle v\rangle$, since $v$ and $\ell(v)$ both have order four. It follows that $X(K)\cap\langle v\rangle=1$. Corollary~\ref{cor:consequences}(e) proves source Example 3.8 for every such chiral facet. No noncanonical or negative-parameter extension of the source's chiral-map assertion is implicit. The quotient $C_4$ used by the proof is a regular marked group; we do not need the one-cell square quotient $\{4,4\}_{(1,0)}$ to be polytopal.

\subsection{The five atomic tight families}
The older input here is Cunningham--Pellicer, \emph{Tight chiral polytopes}, Journal of Algebraic Combinatorics 54 (2021), 837--878. We retain the existence and atomicity of its five Table 4 families having both ends chiral, together with its published Proposition 11. The acquired arXiv version 2009.04566v1 labels the relevant result Proposition 4.2, on printed page 20. The earlier review read that proposition, proof, neighboring embedded-group arguments and the five table rows; it did not independently reprove the full classification.

For those atomic polytopes the retained result gives
\[
              X(P)=\langle\sigma_2^{q'}\rangle,
                     \qquad q/q'\text{ prime},
\]
where $q$ is the middle generator order. The local chirality groups are nontrivial because both ends are chiral. Each embeds in the global $X(P)$: its first-coordinate kernel is obtained by restricting the global mirror mix to an end subgroup. A nontrivial subgroup of the cyclic prime-order $X(P)$ must equal it. Thus
\[
                         X(K)=X(L)=X(P)
\]
as subgroups of the same $\Gamma$, not merely as abstract isomorphism types.

\begin{center}\small
\begin{tabular}{@{}p{.32\textwidth}p{.34\textwidth}p{.24\textwidth}@{}}\toprule
Type & Parameters in the retained family & Common nonidentity element\\\midrule
$\{2m,m^\alpha,2m\}$ & $m$ odd prime, $\alpha\geq2$, $1\leq k\leq m-1$ & $\sigma_2^{m^{\alpha-1}}$, order $m$\\[3pt]
$\{8,2^\beta,8\}$ & $\beta\geq5$, $\epsilon_1,\epsilon_2\in\{\pm1\}$ & $\sigma_2^{2^{\beta-1}}$, order $2$\\[3pt]
$\{2^{\beta-1},2^\beta,2^{\beta-1}\}$ & same range & same element\\[3pt]
$\{8,2^\beta,2^{\beta-1}\}$ & same range & same element\\[3pt]
$\{2^{\beta-1},2^\beta,8\}$ & same range & same element\\\bottomrule
\end{tabular}
\end{center}
The parameters specify the classified marked groups, not every polytope with the same Schläfli type. The indicated elements are nonidentity because the middle generator has exact order $m^\alpha$ or $2^\beta$. Corollary~\ref{cor:consequences}(c) therefore shows that the smallest regular cover of every polytope in these five families is not polytopal. This is the full all-parameter deduction in source Example 3.5. It uses the older classification as stated; finite sampling cannot replace it.
